\documentclass{article}
\usepackage{amsmath,graphicx,algorithm2e,algorithmic,float}
\title{A Study on Computational Metrics}
\author{Bodike Chaithali \\ B24CS013}
\date{\today}
\begin{document}
	\maketitle
	\begin{abstract}
		This report demonstrates essential LATEX features including structured
		sections, equations, tables, figures, algorithms, custom macros, and bibliog-
		raphy management. We illustrate their usage through a new context with
		different formulas.
	\end{abstract}
	\section{Introduction}
	Mathematical relationships are fundamental in computational studies. For exam-
	ple, the Gaussian integral is given by
	\begin{equation}
		\int_{-\infty}^{\infty} e^{-x^2} dx = \sqrt{\pi},
	\end{equation}
	which we will reference in later discussions.
	\section{Methods}
	We measured the refractive index of the material using the custom macro
	\[
		n = \frac{c}{v},
	\]
	where $c$ is the speed of light in vacuum and $v$ is the speed of light in the material.
	\subsection{Figures}
	Figure \ref{fg:1} shows a demonstration image included to illustrate figure handling.
	\begin{figure}[H]
		\centering
		\includegraphics[width=0.4\textwidth]{/home/bodike/Desktop/Sem3/CSP203/Practice/iit-bhilai.jpg}
		\caption{Sample illustration figure.}
		\label{fg:1}
	\end{figure}
	\subsection{Tables}
	Table \ref{tb:1} shows sample measurements of sensor readings.
	\begin{table}[H]
		\centering
		\begin{tabular}{|c|c|c|c|}
			\hline
			Sensor & Reading 1 & Reading 2 & Reading 3 \\
			\hline
			A & 12.5 & 13.1 & 12.9 \\ \hline
			B & 14.0 & 13.8 & 14.2 \\ \hline
			C & 11.8 & 12.0 & 11.9 \\ \hline 
		\end{tabular}
		\caption{Sample sensor readings.}
		\label{tb:1}
	\end{table}
	\section{Algorithm}
	Algorithm \ref{ag:1} demonstrates computing the factorial of a number.\\
	
	\begin{algorithm}[H]
		$n \gets \text{input integer};$ \\
		$result \gets 1;$ \\
		\textbf{for } $i \gets 1$ \textbf{ to } $n$ \textbf{ do} \\
		\hspace*{1em}$result \gets result \times i;$ \\
		\textbf{end} \\
		\caption{Compute factorial of an integer}
		\label{ag:1}
	\end{algorithm}
	
\end{document}